\chapter{Máscaras tiempo-frecuencia}\label{cap:mascaras}

Seleccionar en Fourierizer es construir una función $M(t,f)\in[0,1]$ que dice qué tanto de cada
punto del plano se conserva: $1$ es todo, $0$ es nada y los valores intermedios atenúan. La capa se
resintetiza a partir de la tabla modificada $M\odot\chi$, con la inversa del \cref{cap:stft}:
\[
  y=\mathcal A^{\dagger}\bigl(M\odot\chi\bigr),\qquad (M\odot\chi)(m,k)=M(t_m,f_k)\,\chi(m,k).
\]
Este capítulo describe cómo se construye $M$ a partir de lo que dibujas: rectángulos, lazos,
peines armónicos y la varita. Todo está pensado para dos cosas: que la selección no dependa de la
malla y que suene limpia.

\section{El plano de trabajo: tiempo y octavas}

Las selecciones se definen en el plano $(t,\varphi)$ con $\varphi=\log_2 f$. Una distancia vertical
fija en este plano es un intervalo musical fijo: subir $\varphi$ en $1$ es subir una octava, y
$1/12$ es un semitono. La razón es perceptiva (\cref{cap:sonido}): oímos las alturas por cocientes,
no por diferencias, así que un borde ``suave de un semitono'' significa lo mismo en los graves que
en los agudos. El bin de frecuencia cero, que en esta escala estaría en $-\infty$, se coloca en
$\varphi=\log_2 1=0$.

Una máscara se \emph{rasteriza} evaluándola en los puntos de la malla, $M(m,k)=M(t_m,\log_2 f_k)$.
Como la figura vive en el plano y no en la tabla, la misma selección sirve para cualquier ventana,
tamaño de FFT o salto.

\section{Cada figura es una distancia}

Cada figura $s$ se describe por su \emph{distancia con signo} $d_s(t,\varphi)$ a su borde:
positiva adentro, negativa afuera, cero en el borde. Las distancias se miden en un plano
\emph{escalado} por los dos \emph{feathers} de la figura: $\tau$ en segundos y $\psi$ en octavas,
\[
  (t,\varphi)\ \longmapsto\ \bigl(t/\tau,\ \varphi/\psi\bigr).
\]
Así, una unidad de distancia es ``un feather'': $\tau$ segundos en horizontal o $\psi$ octavas en
vertical. La máscara de la figura es
\[
  M_s=a_s\,\rho(d_s)\,E_s,\qquad
  \rho(d)=\begin{cases}0,&d\le-\tfrac12,\\[2pt]\tfrac12+\tfrac12\sin(\pi d),&|d|<\tfrac12,\\[2pt]1,&d\ge\tfrac12,\end{cases}
\]
donde $a_s\in[0,1]$ es la ganancia de la figura y $E_s$ son sus fades de tiempo
(\cref{sec:fades}).

\paragraph{Rectángulo.} Para $[t_0,t_1]\times[\varphi_0,\varphi_1]$ sean
$d_t=\min(t-t_0,\,t_1-t)/\tau$ y $d_\varphi=\min(\varphi-\varphi_0,\,\varphi_1-\varphi)/\psi$. Entonces
\[
  d=\begin{cases}\min(d_t,d_\varphi), & d_t\ge0\ \text{y}\ d_\varphi\ge0,\\[2pt]
  -\sqrt{\min(d_t,0)^2+\min(d_\varphi,0)^2}, & \text{en otro caso.}\end{cases}
\]

\begin{proposicion}
Esta fórmula es la distancia euclidiana con signo al rectángulo en el plano escalado.
\end{proposicion}

\begin{proof}
Adentro, el punto más cercano del borde está sobre el lado más cercano, a distancia
$\min(d_t,d_\varphi)$. Afuera hay dos casos. Si el punto está frente a un lado (una sola de
$d_t,d_\varphi$ es negativa), la distancia es la de ese lado, $|\min(\cdot,0)|$. Si está frente a una
esquina (ambas negativas), el punto más cercano es la esquina, y por Pitágoras la distancia es
$\sqrt{d_t^2+d_\varphi^2}$. Las dos situaciones son la expresión con $\min(\cdot,0)$.
\end{proof}

\paragraph{Lazo.} Un polígono con vértices $(t_i,\varphi_i)$ usa la distancia mínima a sus segmentos,
calculada en el plano escalado. El signo lo da la \emph{regla par-impar}: un rayo horizontal desde
el punto cruza el contorno un número impar de veces si y solo si el punto está adentro. La regla
funciona igual para lazos cóncavos o que se cruzan a sí mismos.

\section{El borde suave}

\begin{proposicion}[Propiedades de la rampa]\label{prop:rampa}
La rampa $\rho$ es no decreciente, de clase $C^1$, vale $\tfrac12$ en el borde y es
\emph{complementaria}: $\rho(d)+\rho(-d)=1$ para todo $d$.
\end{proposicion}

\begin{proof}
En $|d|<\tfrac12$, $\rho'(d)=\tfrac\pi2\cos(\pi d)\ge0$, y en $d=\pm\tfrac12$ esa derivada vale $0$,
igual que la de los tramos constantes: $\rho$ y $\rho'$ son continuas. La complementariedad es
$\tfrac12+\tfrac12\sin(\pi d)+\tfrac12+\tfrac12\sin(-\pi d)=1$, y en los tramos constantes es
$1+0$.
\end{proof}

La complementariedad tiene una consecuencia muy útil: \textbf{dos figuras que comparten un borde,
con el mismo feather, se reparten la transición sin huecos ni picos}. En cada punto de la franja
común, una aporta $\rho(d)$ y la otra $\rho(-d)$, y suman $1$. Es un \emph{crossfade} espacial: se
puede partir un sonido en regiones y repartirlas entre capas sin que se oiga la costura. Que la
rampa sea $C^1$, y no solo continua, evita el escalón en la pendiente que una rampa lineal tendría
en sus extremos: el borde de la selección no ``chasquea'' en la reconstrucción.

\section{Combinar figuras}

Las figuras se aplican en orden, como capas en un editor de imágenes, empezando con $M=0$:
\[
  \text{sumar:}\quad M\leftarrow M+(1-M)\,M_s,\qquad\qquad
  \text{restar:}\quad M\leftarrow M\,(1-M_s).
\]

\begin{proposicion}\label{prop:union}
Ambas operaciones mantienen $M\in[0,1]$. Si solo hay figuras que suman, el resultado no depende
del orden:
\[
  1-M=\prod_s\bigl(1-M_s\bigr).
\]
\end{proposicion}

\begin{proof}
Si $M,M_s\in[0,1]$, entonces $M+(1-M)M_s$ es un promedio ponderado de $M$ y $1$, y $M(1-M_s)$ es
un producto de números en $[0,1]$. Para el orden, sumar una figura transforma $1-M$ en
$(1-M)-(1-M)M_s=(1-M)(1-M_s)$, así que después de todas las sumas $1-M$ es el producto de los
factores $(1-M_s)$, que conmuta.
\end{proof}

La suma es una ``unión suave'': donde las figuras no se tocan es su unión, y donde se traslapan
no se pasa de $1$. Si $M_s$ fuera la probabilidad de que un punto pertenezca a la figura $s$,
sería la probabilidad de pertenecer a alguna. La resta es una ``diferencia suave''. Una figura
que resta solo afecta a lo que ya estaba seleccionado, y lo que se suma después ya no se ve
afectado: \textbf{el orden importa} entre sumas y restas, igual que con la goma de borrar.

\begin{estudio}
\textbf{Construir una selección.} Empieza con figuras grandes que suman y talla con \emph{Erase}
(o con Alt) lo que sobre. Para hacer un hueco que no se llene, la resta tiene que ir después de
las figuras que cubren esa zona. Para repartir un sonido entre dos capas sin costura, usa figuras
que compartan el borde con el mismo feather.
\end{estudio}

\section{Fades de tiempo}\label{sec:fades}

El feather suaviza el borde de una figura. Los \emph{fades} hacen otra cosa: dan forma al sonido
\emph{dentro} de la selección, como una envolvente. En cada renglón de frecuencia $\varphi$, una
figura cubre uno o varios intervalos de tiempo $[a_j,b_j]$: uno solo en un rectángulo, y los que
corten el contorno en un lazo. La envolvente es
\[
  E(t,\varphi)=\max_j\ S\!\Bigl(\frac{t-a_j}{T_{\text{in}}}\Bigr)\,S\!\Bigl(\frac{b_j-t}{T_{\text{out}}}\Bigr),
  \qquad S(u)=\begin{cases}0,&u\le0,\\ \tfrac12-\tfrac12\cos(\pi u),&0<u<1,\\ 1,&u\ge1,\end{cases}
\]
con la convención de que un fade de duración cero aporta un factor $1$. Cada renglón entra con un
coseno alzado de $T_{\text{in}}$ segundos desde \emph{su propio} borde izquierdo y sale en
$T_{\text{out}}$ antes de su borde derecho. En un lazo inclinado, los fades siguen la inclinación:
un barrido que sube en frecuencia también entra gradualmente a lo largo de su diagonal.

En los renglones que quedan fuera del rango de frecuencias del lazo se usan los intervalos del
renglón cubierto más cercano. Así, la envolvente sigue siendo continua a través del feather de
arriba y de abajo.

\begin{estudio}
\textbf{Feather contra fade.} El feather es corto (milisegundos) y sirve para que un corte no
chasquee. El fade es largo (de cientos de milisegundos a segundos) y es una decisión musical: que
una textura aparezca poco a poco, o que una cola desaparezca en vez de cortarse. Para quitar el
ataque de un sonido y quedarte con su cuerpo, usa un rectángulo con \emph{Fade in} de unos
$50$--$100$ ms.
\end{estudio}

\section{El peine armónico}

Un sonido con altura definida tiene sus parciales cerca de los múltiplos de su fundamental,
$f_0,2f_0,3f_0,\dots$ (\cref{cap:sonido}). El \emph{peine} de Fourierizer selecciona, durante
$[t_0,t_1]$, franjas de ancho $w$ (en octavas) centradas en los armónicos $k f_0$ con $k$ en un
conjunto $K$: los primeros armónicos, solo los impares o solo los pares, hasta $24$ kHz. Su
distancia en frecuencia es la de la franja más cercana,
\[
  d_\varphi=\frac{1}{\psi}\Bigl(\frac w2-\min_{k\in K}\bigl|\varphi-\log_2(kf_0)\bigr|\Bigr),
\]
y se combina con la distancia en tiempo exactamente como en el rectángulo. Así, el feather y los
fades se comportan igual que en cualquier otra figura.

En escala logarítmica, los armónicos se van juntando: entre $k f_0$ y $(k+1)f_0$ hay
$\log_2\frac{k+1}{k}$ octavas, que tiende a cero.

\begin{proposicion}
Las franjas de los armónicos $k$ y $k+1$ se tocan si y solo si
\[
  k\ \ge\ \frac{1}{2^{w}-1}.
\]
A partir de ahí, el peine ya no separa armónicos: selecciona una banda continua.
\end{proposicion}

\begin{proof}
Cada franja se extiende $w/2$ a cada lado de su centro, así que dos franjas vecinas se tocan cuando
la separación de sus centros es a lo más $w$: $\log_2\frac{k+1}{k}\le w$, es decir,
$1+\frac1k\le 2^w$, que equivale a la desigualdad del enunciado.
\end{proof}

Con un ancho de 100 cents ($w=1/12$), las franjas se juntan desde el armónico 17. Con 50 cents,
desde el 35, y con 200 cents, desde el 9. Esto coincide con lo que se oye: los armónicos altos de
una nota grave están más cerca entre sí que un semitono, y un oído (o una selección) ya no los
separa.

\begin{estudio}
\textbf{Usar el peine.}
\begin{itemize}
  \item \emph{Aislar una nota en una mezcla:} un peine sobre su fundamental con 50--100 cents de
        ancho toma la nota y deja casi todo lo demás. Funciona mejor con ventanas largas
        (\cref{cap:stft}).
  \item \emph{Odd} deja solo los armónicos impares. El timbre se vuelve hueco y nasal, como un
        clarinete o una onda cuadrada. \emph{Even} deja solo los pares: suena una octava arriba y
        más brillante.
  \item Si la nota tiene vibrato o glissando, la fundamental se mueve y un peine fijo la pierde por
        momentos. Sube el ancho, o usa varios peines cortos.
\end{itemize}
\end{estudio}

\section{La varita}

La varita no se define por una fórmula sino por el propio sonido. Al dar clic en un punto del
espectrograma, se seleccionan los píxeles \emph{conectados} a ese punto cuyo nivel está a menos de
una tolerancia $\delta$ (en dB) por debajo del nivel del punto. El borde de esa región se traza a lo
largo de las aristas de los píxeles y se simplifica con el algoritmo de Ramer–Douglas–Peucker
hasta un máximo de 120 vértices. El resultado es un lazo ordinario, de modo que todo lo anterior
(distancia, feather, fades) se le aplica sin cambios.

La varita depende de la imagen que ves: de la ventana y del rango de la vista. El lazo que produce,
en cambio, vive en el plano $(t,\varphi)$ como cualquier otro, y no cambia si después modificas la
malla.

\section{Artefactos y cómo evitarlos}

Una máscara modifica la tabla, y la reconstrucción proyecta el resultado sobre las tablas
consistentes (\cref{prop:proyeccion}). De ahí salen los tres artefactos clásicos del filtrado
tiempo-frecuencia:
\begin{itemize}
  \item \textbf{Ringing.} Un borde duro no es consistente. La proyección lo redondea y puede dejar
        oscilaciones cerca del borde, sobre todo en frecuencia: tonos que ``suenan'' un poco fuera
        de la selección. \emph{Remedio:} feather.
  \item \textbf{Ruido musical.} Una máscara que salta entre $0$ y $1$ de coeficiente en coeficiente
        (por ejemplo, una varita sobre ruido, con tolerancia baja) deja islas sueltas de energía.
        Cada isla suena como un tono breve y aleatorio. \emph{Remedio:} figuras más grandes y
        lisas, o más tolerancia.
  \item \textbf{Pre-eco.} Una selección que empieza justo en un ataque deja pasar la parte de las
        ventanas que lo contienen y que empieza antes, del orden de $L/2$. El golpe aparece un poco
        antes, y suavizado. \emph{Remedio:} ventanas cortas en percusiones, o multirresolución
        (\cref{cap:multirresolucion}).
\end{itemize}
